\begin{lemma}
Let \(P\) be a finite set of length-\(k\) paths, and let
\(\sigma:P\to\{0,1\}^k\) be a signature map, where a \(1\) records an
unmarked step and a \(0\) records a marked step.  Let \(w\le k\) and
\(h\le\Delta\).  Suppose that every \(p\in P\) has at most \(w\) unmarked
steps, and suppose that for every binary string \(z\in\{0,1\}^k\), the number
of paths with signature \(z\) is at most
\[
  \Delta^{|z|}h^{k-|z|},
\]
where \(|z|\) is the number of \(1\)'s in \(z\).  Then
\[
  |P|\le 2^k\Delta^w h^{k-w}.
\]
\end{lemma}
\begin{proof}
For a binary string \(z\), write \(P_z=\{p\in P:\sigma(p)=z\}\).  The sets
\(P_z\), as \(z\) ranges over \(\{0,1\}^k\), partition \(P\).  Hence
\[
  |P|=\sum_{z\in\{0,1\}^k}|P_z|.
\]
If \(P_z\neq\emptyset\), choose \(p\in P_z\).  Then the hypothesis on paths
gives \(|z|=|\sigma(p)|\le w\).  If \(P_z=\emptyset\), its contribution is
zero and the same upper bound may be used harmlessly.  Thus only signatures
with \(|z|\le w\) can contribute.

For each contributing \(z\), the fiber hypothesis gives
\[
  |P_z|\le \Delta^{|z|}h^{k-|z|}. \tag{1}
\]
Because \(h\le\Delta\) and \(|z|\le w\le k\), the right hand side in (1) is
at most \(\Delta^w h^{k-w}\).  Indeed, if \(h=0\), then either \(k=w\), in
which case this is immediate from \(|z|\le w=k\), or \(k>w\), in which case
every nonzero term has \(k-|z|>0\) and is already \(0\).  If \(h>0\), moving
one unit of exponent from \(h\) to \(\Delta\) multiplies the monomial by
\(\Delta/h\ge1\), and repeating this \(w-|z|\) times gives the claim.

There are exactly \(2^k\) binary strings of length \(k\).  Summing the
uniform bound \(\Delta^w h^{k-w}\) over all signatures therefore yields
\[
  |P|
  \le \sum_{z\in\{0,1\}^k}\Delta^w h^{k-w}
  =2^k\Delta^w h^{k-w}.
\]
\end{proof}
